% ============================================================================
% MUSTERLÖSUNG: Reibung Vertiefung – ABS & Anwendungen (UE 6)
% Physik Klasse 7 - Gesamtschule MV
% ============================================================================
\documentclass[a4paper,11pt]{article}

\usepackage[left=1.5cm,right=1.5cm,top=1.5cm,bottom=1.5cm]{geometry}
\usepackage[utf8]{inputenc}
\usepackage[T1]{fontenc}
\usepackage[ngerman]{babel}
\usepackage{tikz}
\usepackage{fancyhdr}
\usepackage{tcolorbox}
\usepackage{amssymb,amsmath}
\usepackage{xcolor}
\usepackage{tabularx}
\usepackage{colortbl}

\renewcommand{\familydefault}{\sfdefault}

% Seiteneinstellungen
\pagestyle{fancy}
\fancyhf{}
\renewcommand{\headrulewidth}{0.4pt}
\lhead{\textbf{Physik Klasse 7}}
\chead{Reibung Vertiefung -- \textcolor{red}{MUSTERLÖSUNG}}
\rhead{}
\setlength{\headheight}{14pt}
\cfoot{}

% Farben
\definecolor{physik}{HTML}{2c5aa0}
\definecolor{success}{HTML}{2a7a4b}
\definecolor{warning}{HTML}{b35c00}
\definecolor{loesung}{HTML}{c62828}

% Hilfsbefehle
\newcommand{\luecke}[1]{\underline{\hspace{#1}}}
\newcommand{\stern}[1]{\textcolor{warning}{\textbf{#1}}}
\newcommand{\lsg}[1]{\textcolor{loesung}{\textbf{#1}}}

\setlength{\parindent}{0pt}
\setlength{\parskip}{4pt}

\begin{document}

% ==================== KOPFBEREICH ====================

\begin{tcolorbox}[colback=white, colframe=physik, boxrule=1.5pt, arc=0pt,
    title={\textbf{Reibung Vertiefung: ABS \& Anwendungen -- \textcolor{red}{MUSTERLÖSUNG}}}, fonttitle=\large, coltitle=black, colbacktitle=white]

\begin{tabular}{@{}ll@{\hspace{2cm}}ll@{}}
\textbf{Klasse:} & \lsg{7A/7B} & \textbf{Datum:} & \lsg{\_\_\_\_\_\_\_\_\_\_}
\end{tabular}

\end{tcolorbox}

\vspace{2mm}

% ==================== K1: Haftreibungsgrenze ====================

\begin{tcolorbox}[colback=white, colframe=physik, boxrule=1pt, arc=0pt,
    title={\textbf{K1: Die Haftreibungsgrenze}}, fonttitle=\normalsize, coltitle=black, colbacktitle=white]
Die \textbf{Haftreibungsgrenze} ist die maximale Kraft, die die Haftreibung aufbringen kann.\\[1mm]
Wird sie überschritten, beginnt der Körper zu \lsg{gleiten} (gleiten/haften).\\[1mm]
\textbf{Formel:} $F_{\text{Haft,max}} = \mu_{\text{Haft}} \cdot F_N$
\end{tcolorbox}

\vspace{2mm}

% ==================== B1: Beobachtung Haftreibungsgrenze ====================

\textbf{B1: Beobachtung – Wann reißt die Haftreibung ab?} \hfill (2 BE)

Beobachte die Simulation: Bei welcher Zugkraft beginnt der Klotz zu gleiten?

$F_{\text{Haft,max}} = $ \lsg{8} N

Was passiert mit der Reibungskraft, wenn der Klotz gleitet? \lsg{Sie wird kleiner (Gleitreibung)}

\vspace{2mm}

% ==================== K2: ABS ====================

\begin{tcolorbox}[colback=white, colframe=physik, boxrule=1pt, arc=0pt,
    title={\textbf{K2: So funktioniert ABS}}, fonttitle=\normalsize, coltitle=black, colbacktitle=white]
\textbf{ABS} = Antiblockiersystem\\[1mm]
\begin{tabular}{@{}ll@{}}
Blockiertes Rad: & \textcolor{red}{Gleitreibung} ($\mu \approx 0{,}5$) $\rightarrow$ längerer Bremsweg \\
Rollendes Rad: & \textcolor{success}{Haftreibung} ($\mu \approx 0{,}8$) $\rightarrow$ kürzerer Bremsweg \\
\end{tabular}\\[1mm]
ABS hält die Räder knapp unter der \lsg{Haftreibungsgrenze}.
\end{tcolorbox}

\vspace{2mm}

% ==================== T1: ABS-Vergleich ====================

\textbf{T1: Bremsweg-Vergleich mit und ohne ABS} \hfill (2 BE)

\begin{center}
\begin{tabular}{|l|c|c|c|}
\hline
\rowcolor{white}
 & \textbf{Reibungsart} & \textbf{Bremskraft} & \textbf{Bremsweg} \\
\hline
Ohne ABS (blockiert) & Gleitreibung & \lsg{500} N & \lsg{10,0} m \\
\hline
Mit ABS (rollt) & Haftreibung & \lsg{800} N & \lsg{6,3} m \\
\hline
\end{tabular}
\end{center}

Ersparnis durch ABS: Der Bremsweg ist ca. \lsg{37} \% kürzer.

\vspace{2mm}

% ==================== B2: Rollreibung Vergleich ====================

\textbf{B2: Rollreibung – Auto vs. Eisenbahn} \hfill (2 BE)

\begin{tabular}{@{}p{7cm}p{7cm}@{}}
\textbf{Auto (Gummi/Asphalt):} $\mu_{\text{Roll}} = 0{,}015$ & \textbf{Zug (Stahl/Stahl):} $\mu_{\text{Roll}} = 0{,}001$ \\
\end{tabular}

Welches Fahrzeug hat die größere Rollreibung? \lsg{Auto}

Nenne einen \textbf{Vorteil} der geringen Eisenbahn-Reibung: \lsg{Energieeffizienter / weniger Kraftstoff}

Nenne einen \textbf{Nachteil} der geringen Eisenbahn-Reibung: \lsg{Längerer Bremsweg}

\vspace{3mm}

% ==================== R1: Rechnung ====================

\begin{tcolorbox}[colback=white, colframe=warning, boxrule=1pt, arc=0pt,
    title={\textbf{R1: Rechnung \stern{**}}}, fonttitle=\normalsize, coltitle=black, colbacktitle=white]
Ein Schrank wiegt 800 N und steht auf Holzboden ($\mu = 0{,}5$). Berechne die Reibungskraft.

\textbf{Gegeben:} $F_N = $ \lsg{800 N}, $\mu = $ \lsg{0,5}

\textbf{Formel:} $F_R = $ \lsg{$\mu \cdot F_N$}

\textbf{Einsetzen:} $F_R = $ \lsg{$0{,}5 \cdot 800\,\text{N}$}

\textbf{Ergebnis:} $F_R = $ \lsg{400} N
\end{tcolorbox}

\vspace{2mm}

% ==================== R2: Umstellen ====================

\begin{tcolorbox}[colback=white, colframe=warning, boxrule=1pt, arc=0pt,
    title={\textbf{R2: Umstellen \stern{***}}}, fonttitle=\normalsize, coltitle=black, colbacktitle=white]
Ein Auto bremst bei Regen ($\mu = 0{,}5$). Es soll dieselbe Bremskraft erreichen wie bei trockenem Asphalt: $F_R = 9600$ N.

Welche Normalkraft $F_N$ bräuchte das Auto?

\textbf{Umgestellte Formel:} $F_N = $ \lsg{$F_R / \mu$}

\textbf{Einsetzen:} $F_N = $ \lsg{$9600\,\text{N} / 0{,}5$}

\textbf{Ergebnis:} $F_N = $ \lsg{19.200} N

\textit{Bedeutung: Das Auto müsste fast \lsg{1,6} mal so schwer sein!}
\end{tcolorbox}

\vspace{3mm}

% ==================== Punktekasten ====================

\begin{center}
\begin{tabular}{|c|c|c|}
\hline
\rowcolor{white}
\textbf{\stern{*} Basis} & \textbf{\stern{**} Standard} & \textbf{\stern{***} Erweitert} \\
\hline
\lsg{8} / 8 BE & \lsg{12} / 12 BE & \lsg{16} / 16 BE \\
\hline
\end{tabular}
\end{center}

\end{document}
